% \section{Related Work and Existing Tools}
% \label{sec:related-work-and-existing-tools}

% Few existing tools directly address the combination of 
% method selection, knowledge representation, and 
% implementation support that this thesis investigates. 
% Most related tools either focus on automating a specific 
% part of the ML workflow or provide broad access to AI 
% resources without guiding users through choosing a 
% suitable method for their problem. The following sections 
% review the most relevant examples found.

% \subsection{Ontology-Based Meta AutoML}
% \label{sec:ontology-based-meta-automl}

% OMA-ML \parencite{zender2022ontology_based_meta_automl} 
% is a system that combines multiple AutoML solutions 
% through the ML Ontology, described in Section~\ref{sec:ml-ontology-structure-and-content}.
% The ML Ontology is used to filter configuration 
% options and select appropriate AutoML strategies, 
% making the process more transparent than standard 
% AutoML systems. It also supports a wider range of 
% methods by integrating solutions across multiple 
% ML libraries rather than being tied to a single one.
% Like other AutoML tools, OMA-ML requires the user 
% to provide a labelled dataset and define the ML task 
% upfront. 

% \subsection{Template-Based Code Generation}
% \label{sec:template-based-code-generation}

% Traingenerator \parencite{rieke2021traingenerator} 
% is a web-based tool that generates template Python 
% code and Jupyter notebooks for ML tasks. Users select 
% a task and framework from a menu and receive a working 
% code template covering common steps such as data 
% loading, model setup, training, and evaluation. 
% The generated code can be downloaded as a Python 
% script or notebook, and optionally opened directly 
% in Google Colab.

% The tool assumes the user has already decided on a 
% method and framework before use. As of May 2026,
% the hosted version of Traingenerator 
% produces a runtime error and does not render output 
% correctly. The tool appeared to function as intended 
% during the development period of this thesis. 
% The tool can still be inspected through the GitHub 
% repository.\footnote{\url{https://github.com/jrieke/traingenerator}}

% \subsection{AI-on-Demand Platform}
% \label{sec:ai-on-demand-platform}
% The European AI-on-Demand Platform (AIoDP) and its 
% DeployAI implementation offer a broad marketplace of 
% pre-trained models, reusable AI modules, and 
% domain-specific applications supported by cloud 
% and HPC infrastructure 
% \parencite{troumpoukis2025deployai_ai_on_demand_platform}. 
% The platform targets SMEs, researchers, and public 
% sector organisations and provides tools for building, 
% training, and deploying AI solutions at scale. 
% Users can browse and filter available solutions 
% to narrow down options relevant to their domain 
% or use case. The primary focus is on access and 
% deployment of existing solutions rather than on 
% guiding users through method selection for a 
% new problem.

% \subsection{Approaches to Method Selection}
% \label{sec:approaches-to-method-selection}

% Selecting an appropriate machine learning method is not straightforward. 
% As discussed in Section~\ref{sec:motivation}, the large number of
% available algorithms, each with distinct strengths and limitations, creates 
% a real challenge for practitioners. This challenge is formally described by 
% the Algorithm Selection Problem \parencite{smith_miles2009algorithm_selection}, 
% which asks which algorithm is likely to perform best for a given problem. 
% The No Free Lunch theorems provide a theoretical basis for this difficulty, 
% showing that no single algorithm performs best across all problems and that 
% good selection requires understanding the characteristics of the problem at 
% hand \parencite{wolpert1997no_free_lunch}.

% Several practical tools have been developed to help with method selection. 
% Scikit-learn provides a publicly available decision flowchart that guides 
% users through questions about data size and task type to suggest candidate 
% algorithms \parencite{sklearn_algorithm_cheatsheet}, and Microsoft Azure 
% offers a similar cheat sheet for its machine learning platform 
% \parencite{azure_algorithm_cheatsheet}. While these tools are accessible 
% and widely used, they have clear limitations. 
% The selection criteria used in both tools are mainly accuracy and prediction 
% speed, something that can in practice lead to poor results \parencite{sala2018selecting_ml_algorithm}. 
% Factors such as interpretability, data characteristics, and domain context 
% are not considered, nor do they offer any guidance on how the problem 
% itself should be formulated \parencite{chen2023ml_manufacturing_industry4}.

% More structured frameworks have been proposed in the academic literature. 
% \textcite{sala2018selecting_ml_algorithm} present a two-layer selection 
% framework that incorporates algorithm characteristics such as scalability 
% and robustness to outliers, in addition to learning type and response type. 
% Domain-specific frameworks have also been proposed, such as that of 
% \textcite{arena2024conceptual_framework_ml_algorithm_selection_predictive_maintenance}, 
% who develop a structured set of decision variables for algorithm selection 
% in predictive maintenance. The authors note that the algorithm is often chosen by comparing 
% candidates only on predictive and computational performance, which 
% can lead to results that are poor or too generic for the problem at 
% hand, and that no tool relates algorithmic choices to the 
% application domain.

% A related class of approaches is automated machine learning (AutoML), 
% which seeks to automatically select, compose, and parametrise machine 
% learning models to achieve optimal performance on a given task or dataset 
% \parencite{waring2020automated_machine_learning_review_state_of_the_art_opportunities_healthcare}. 
% AutoML systems automate parts of the machine learning pipeline, including 
% data preparation, feature engineering, hyperparameter optimisation, and 
% model generation, with the goal of reducing the demand for expert knowledge 
% and making machine learning more accessible to domain experts 
% \parencite{he2021automl_survey_state_of_the_art}. 
% Within AutoML, the algorithm selection problem is typically treated as an
% optimisation task, where dataset characteristics are used to predict which 
% algorithm is likely to perform best 
% \parencite{barbudo2023eight_years_automl_categorisation_review_trends}.

% However, AutoML approaches have notable limitations in the context of 
% decision support for non-experts. Most systems treat the selection and 
% configuration process as a black box, providing little explanation of why 
% a particular configuration was chosen 
% \parencite{barbudo2023eight_years_automl_categorisation_review_trends}. 
% This lack of transparency is a significant barrier to adoption, as users 
% who do not understand or trust the output of a system are unlikely to apply 
% it in practice 
% \parencite{waring2020automated_machine_learning_review_state_of_the_art_opportunities_healthcare}. 
% Furthermore, AutoML primarily addresses the later stages of the machine 
% learning pipeline, and the phases considered inherently human, such as 
% problem formulation and interpretation of results, have received little 
% attention \parencite{barbudo2023eight_years_automl_categorisation_review_trends}. 
% This means that AutoML does not address the earlier challenge of helping 
% practitioners understand which type of problem they have and which general 
% approach is appropriate. This gap motivates the 
% work in this thesis.

\section{Related Work and Existing Tools}
\label{sec:related-work-and-existing-tools}

Few existing tools or approaches address the combination of 
method selection, knowledge representation, and implementation 
support that this thesis investigates. The work reviewed here 
falls into two broad groups. The first concerns method 
selection: tools, frameworks, and knowledge representations 
that help a user identify a suitable machine learning method 
for a problem. The second concerns semantic retrieval from 
scientific literature, which is central to how method selection 
is approached in this thesis. The two groups are reviewed in 
turn, followed by a summary of the gaps that motivate the work.

\subsection{Tools and Platforms}
\label{sec:tools-and-platforms}

Several tools provide access to machine learning methods or 
automate parts of the workflow without guiding the user through 
the choice of a method.

OMA-ML \parencite{zender2022ontology_based_meta_automl} combines 
multiple AutoML solutions through the ML Ontology, described in 
Section~\ref{sec:ml-ontology-structure-and-content}. The 
ontology is used to filter configuration options and select 
appropriate AutoML strategies, which makes the process more 
transparent than standard AutoML systems, and it supports a 
wider range of methods by integrating solutions across multiple 
machine learning libraries rather than being tied to a single 
one. Like other AutoML tools, OMA-ML requires the user to 
provide a labelled dataset and to define the machine learning 
task in advance.

Traingenerator \parencite{rieke2021traingenerator} is a 
web-based tool that generates template Python code and Jupyter 
notebooks for machine learning tasks. A task and framework are 
selected from a menu, and a working code template is produced 
that covers common steps such as data loading, model setup, 
training, and evaluation. The generated code can be downloaded 
as a script or notebook, or opened directly in Google Colab. 
The tool assumes that the method and framework have already 
been chosen before use. As of May 2026, the hosted version 
produces a runtime error and does not render output correctly, 
although it appeared to function as intended during the 
development period of this thesis, and it can still be inspected 
through the GitHub 
repository.\footnote{\url{https://github.com/jrieke/traingenerator}}

The European AI-on-Demand Platform (AIoDP) and its DeployAI 
implementation offer a broad marketplace of pre-trained models, 
reusable AI modules, and domain-specific applications supported 
by cloud and HPC infrastructure 
\parencite{troumpoukis2025deployai_ai_on_demand_platform}. The 
platform targets SMEs, researchers, and public sector 
organisations, and provides tools for building, training, and 
deploying AI solutions at scale. Available solutions can be 
browsed and filtered to narrow down options relevant to a 
domain or use case. The focus is on access to and deployment of 
existing solutions rather than on guiding the user through 
method selection for a new problem.

\subsection{Approaches to Method Selection}
\label{sec:approaches-to-method-selection}

Selecting an appropriate machine learning method is not 
straightforward. As discussed in Section~\ref{sec:motivation}, 
the large number of available methods, each with distinct 
strengths and limitations, creates a real challenge for 
practitioners. This challenge is formally described by the 
Algorithm Selection Problem 
\parencite{smith_miles2009algorithm_selection}, which asks which 
algorithm is likely to perform best for a given problem. The No 
Free Lunch theorems provide a theoretical basis for the 
difficulty, showing that no single algorithm performs best 
across all problems and that good selection requires 
understanding the characteristics of the problem at hand 
\parencite{wolpert1997no_free_lunch}.

The most accessible tools for method selection are decision 
flowcharts. Scikit-learn provides a publicly available 
flowchart that guides users through questions about data size 
and task type to suggest candidate algorithms 
\parencite{sklearn_algorithm_cheatsheet}, and Microsoft Azure 
offers a similar cheat sheet for its machine learning platform 
\parencite{azure_algorithm_cheatsheet}. While widely used, both 
rely mainly on accuracy and prediction speed as selection 
criteria, which can in practice lead to poor results 
\parencite{sala2018selecting_ml_algorithm}. Factors such as 
interpretability, data characteristics, and domain context are 
not considered, and neither tool offers guidance on how the 
problem itself should be formulated 
\parencite{chen2023ml_manufacturing_industry4}.

More structured frameworks characterise methods against a set 
of criteria and match a described problem to them. 
\textcite{riesener2020} present an evaluation method for 
selecting a method in product development, aimed at users 
without in-depth machine learning knowledge. From a large set 
of methods found in the literature, the eleven most frequently 
described were selected and each rated on nine criteria, 
covering the learning task, the form of learning, and seven 
performance properties such as accuracy, training duration, 
computational effort, and transparency. The result is a matrix 
of method strengths and weaknesses; the user describes a 
problem in the same criteria, and the method with the smallest 
Euclidean distance to that description is suggested. 
\textcite{lickert2021_ml_selection_reverse_logistics} take a 
similar approach for supervised methods, describing seven 
methods against criteria grouped into input, implementation, 
and output, and matching a problem to them through a single 
overview table. The approach deliberately supports a 
preselection of candidate methods for later testing rather than 
the choice of a single optimal method, and is demonstrated on a 
case from reverse logistics. Both frameworks derive method 
properties manually, and the criteria used overlap considerably 
with the performance attributes represented in this thesis.

Other frameworks focus on the structure of the matching process 
itself. \textcite{sala2018selecting_ml_algorithm} propose a 
two-layer selection framework that incorporates method 
characteristics such as scalability and robustness to outliers 
alongside learning type and response type. 
\textcite{arena2024conceptual_framework_ml_algorithm_selection_predictive_maintenance} 
develop a structured set of decision variables for method 
selection in predictive maintenance, noting that methods are 
often chosen by comparing candidates only on predictive and 
computational performance, which can give results that are poor 
or too generic, and that no tool relates the choice to the 
application domain. A more elaborate matching mechanism is 
proposed by \textcite{sonntag2023pattern_ml_product_development}, 
who apply a pattern language to method selection in product 
development. Each method is characterised through a pattern with 
fixed attributes, and the user formalises a problem using the 
same attributes, so that problem and method can be matched 
directly. The concept was demonstrated for two methods and 
tested with seven engineers without machine learning knowledge; 
the correct method was identified in 64\% of cases, and the 
errors arose mainly from incorrect formalisation of the problem 
rather than from the matching itself. The same group later 
proposed a multi-criteria decision analysis framework 
\parencite{sonntag2024_mcda_framework}, in which a reference 
process first maps a task to a class of methods, and the 
remaining methods are then ranked against qualitative criteria 
using the PROMETHEE outranking method, weighted by the user's 
stated preferences.

A related class of approaches is automated machine learning 
(AutoML), which seeks to automatically select, compose, and 
parametrise models to achieve good performance on a given task 
or dataset 
\parencite{waring2020automated_machine_learning_review_state_of_the_art_opportunities_healthcare}. 
AutoML automates parts of the pipeline, including data 
preparation, feature engineering, hyperparameter optimisation, 
and model generation, with the goal of reducing the demand for 
expert knowledge 
\parencite{he2021automl_survey_state_of_the_art}. The selection 
problem is typically treated as an optimisation task, where 
dataset characteristics are used to predict which method is 
likely to perform best 
\parencite{barbudo2023eight_years_automl_categorisation_review_trends}. 
For decision support aimed at non-experts, however, AutoML has 
notable limitations. Most systems treat the selection and 
configuration process as a black box, with little explanation 
of why a particular configuration was chosen 
\parencite{barbudo2023eight_years_automl_categorisation_review_trends}, 
and this lack of transparency is a barrier to adoption, since 
users who do not understand or trust the output are unlikely to 
apply it 
\parencite{waring2020automated_machine_learning_review_state_of_the_art_opportunities_healthcare}. 
AutoML also addresses mainly the later stages of the pipeline, 
while the phases that are inherently human, such as problem 
formulation and interpretation of results, receive little 
attention 
\parencite{barbudo2023eight_years_automl_categorisation_review_trends}. 
It therefore does not address the earlier difficulty of helping 
a practitioner understand which kind of problem they have and 
which general approach is appropriate.

\subsection{Knowledge Representation for Method Selection}
\label{sec:knowledge-representation-for-method-selection}

A separate line of work represents method knowledge in a 
structured, queryable form rather than in a fixed table or 
flowchart. \textcite{gerschutz2021semantic_web_approach} 
present a semantic knowledge base for data-driven methods in 
product development, aimed at SMEs that lack machine learning 
expertise. Method knowledge is structured in a class model in 
which each method is described by properties such as its 
potentials, limits, algorithm, product development use cases, 
and the tool that implements it. The model is realised as a 
demonstrator in Semantic MediaWiki, where the knowledge is 
captured through semantic links and retrieved through queries, 
so that a user can query for a method that supports a given 
activity such as requirements analysis. The method properties 
are populated manually from literature review and industrial 
use cases.

This work was later developed into a formal ontology. AI4PD 
\parencite{gerschutz2023_AI4DP} links data-driven methods to 
product development processes through an OWL ontology, modelled 
in Protégé and queried with SPARQL. It introduces concepts for 
processes, methods, tools, and the people involved, and uses a 
task cluster concept to connect company-specific processes to 
the methods that can support them. A problem is described 
formally, in terms of the task, the available data, and the 
hardware, and a SPARQL query then returns the matching methods 
and tools along with documented use cases. The authors note 
that formulating such a query is difficult for inexperienced 
users and suggest that a graphical interface would help, and 
that the manual effort of populating the ontology with method 
knowledge is high. Both observations are relevant to the design 
of any ontology-based method knowledge base. The ontology 
underlying AI4PD is the same ML Ontology family discussed in 
Section~\ref{sec:ml-ontology-structure-and-content}, which this 
thesis also builds on.

\subsection{Semantic Retrieval from Scientific Literature}
\label{sec:semantic-retrieval-from-scientific-literature}

The second group of related work concerns retrieving relevant 
documents from scientific literature based on the meaning of a 
query rather than on exact keywords. The underlying techniques, 
text embeddings, cosine similarity, and the bi-encoder and 
cross-encoder architectures, are described in 
Section~\ref{sec:semantic-retrieval-and-text-similarity}. This 
section reviews systems that apply these techniques to 
scientific literature.

\textcite{berenguer2024retrieve_data_from_research} retrieve 
data tables from publications based on a research problem stated 
as an abstract. The abstract serves as a representation of the 
research problem, and the semantic similarity between the 
abstract and the indexed tables is computed from text 
embeddings to return a ranked list of relevant tables. The work 
argues that a rich natural-language formulation of a research 
problem gives better retrieval than isolated keywords, and that 
matching such a description against scientific content is 
feasible using embeddings. A comparison of nine language models 
finds that contextual models, and in particular models 
pre-trained on scientific text, clearly outperform 
non-contextual ones, and that surrounding context is important 
for good results. The principle of using a free-text problem 
description as a query, matched against scientific content 
through embedding similarity, is the same principle that 
underlies the retrieval approach in this thesis, although the 
output there is data tables rather than methods.

\textcite{schopf2024nlp} present NLP-KG, a system for 
exploratory search of natural language processing literature. 
It is built around a knowledge graph that links fields of 
study, publications, authors, and venues, with a 
semi-automatically constructed hierarchy of fields of study. 
The structural knowledge graph and the publication embeddings 
are stored separately, the former in a graph database and the 
latter in a vector database, so that the graph handles 
relations and the vector database handles semantic similarity 
search. Retrieval combines sparse and dense representations, 
which are merged using Reciprocal Rank Fusion and then reranked 
using metadata such as citation count. The same separation of a 
structural knowledge store from a vector store for retrieval is 
present in this thesis, and Reciprocal Rank Fusion is used here 
as well, although for combining method-level rankings rather 
than two document retrievers. NLP-KG is intended to help 
researchers explore an unfamiliar field rather than to select a 
method for a problem.

Two further approaches concern retrieval where the query is a 
full paper rather than a problem description. 
\textcite{park2025chain_of_retrieval} retrieve related papers 
by decomposing a query paper into aspect-specific subqueries and 
iteratively expanding the search, and report that relying only 
on abstracts is suboptimal because abstracts offer only sparse, 
high-level summaries of a paper. 
\textcite{ostendorff2022_aspect_based_embeddings} recommend 
related machine learning papers using separate embeddings for 
the task, method, and dataset of each paper, and find that 
generic embeddings of titles and abstracts capture the dataset 
of a paper well but its method poorly, because methodological 
detail is mentioned only marginally in titles and abstracts. 
Both observations are relevant to a system that retrieves 
methods through the abstracts of the articles that mention 
them, and they are returned to in the discussion.

\subsection{Summary of Gaps}
\label{sec:summary-of-gaps}

Across the work reviewed here, several gaps remain. Decision 
flowcharts and AutoML systems support method selection but rely 
mainly on performance criteria, treat the process as a black 
box, or assume that the task and a labelled dataset are already 
defined, and so give little help with the earlier and 
inherently human step of understanding the problem and choosing 
a general approach. The structured frameworks and knowledge 
representations address this step more directly, by 
characterising methods and matching them to a described 
problem, but they depend on method knowledge that is populated 
manually, are demonstrated on a small number of methods or a 
single use case, and require the user to formalise the problem 
in a fixed structure, which has been shown to be a source of 
error. Semantic retrieval from literature offers a way to match 
a free-text problem description against a large body of methods 
without manual formalisation, but existing systems apply it to 
exploratory search or to retrieving tables or related papers, 
rather than to selecting a method and supporting its 
implementation. This thesis addresses these gaps together, and 
the approach is described in the following chapter.